\section{Introduction}

Selecting an algorithm from instance features is a classical problem \cite{rice1976algorithm}.  Modern instance-space analysis turns features and performance into interpretable regions that reveal where algorithms succeed or fail \cite{smithmiles2009meta,smithmiles2023isa}.  The same perspective has reached quantum optimization and compilation.  It has been used to study parameter initialization \cite{katial2025qaoa}, quantum versus classical algorithm choice \cite{moussa2020selection}, compilation-option prediction \cite{quetschlich2023options}, and device-specific compiler selection \cite{quetschlich2025predictor}.  These studies establish that feature-conditioned performance mapping and algorithm selection are prior work, not contributions of this paper.

Exact compiler models permit a different question.  When an exact optimum or a complete proof is available, one can ask why a reference compiler fails, whether a restricted family always contains an optimum, and whether the feature representation itself determines the label.  These questions have different logical answers.  A support bound may be intrinsic or merely sufficient for one proof.  A proof can hold only for a region of objective weights.  A feature map can fit a complete finite domain yet fail on new sizes.  Combining these answers into one accuracy surface hides the distinction.

We therefore use a \emph{regime map} as a typed reporting record rather than a single phase diagram.  The record separates the reference-exact region, constructive improvement mechanisms, intrinsic support, proof-derived support ceilings, objective-conditioned proof validity, and feature determination with prospective transfer.  A negative result is part of the record.  For example, mixed feature cells prove that no classifier using those features can be exact on the stated domain, while a failed prospective forecast shows that an in-domain rule did not transfer.  Missing coordinates are marked as unknown or not applicable rather than silently inferred from another coordinate.

The evidence comes from four frozen quantum-compilation models (Table~\ref{tab:cross-model}).  Two Pauli block-encoding models admit all-size support results.  A six-term linear-combination model admits an exact incumbent boundary.  A weighted Clifford stabilizer-state model supplies complete finite graphs and a prospective size-transfer test.  The models are deliberately heterogeneous.  They are examples that separate logical coordinates, not a statistical sample of compilers.

The primary contributions are the three exact, model-specific results and the retained adverse transfer study.  The regime map is their organizing schema.  Together they establish a concrete comparison rule: numerical support, accuracy or region values may be compared only after their model, objective, domain and evidence class have been aligned.  This changes the interpretation of several results.  A safe syndrome-rank ceiling of five is not evidence that support five is intrinsic when a different proof gives intrinsic support one.  Silence outside a proof-validity region is not evidence of necessity.  Elimination of all finite-domain feature collisions is not evidence of transfer when the next-size test has almost no coverage.  The distinctions therefore prevent specific overclaims that a single performance surface or separately reported headline numbers would leave open.
